\begin{frame}{자주 하는 실수: 시계열을 무작위로 나누기}
  환자에게 있는 ``단위'' 문제의 다른 얼굴이 \textbf{시간 순서}.
  시계열은 \textbf{데이터가 이미 시간순 정렬되어 있을 때 무작위 shuffle이
  구조를 깨는} 방향으로 작용.
  \vskip0.3em
  \small
  추세(매일 $+1$)가 있는 일일 데이터 365일, kNN:
  \begin{tabular}{@{}ll@{}}
    \toprule
    분할 방식 & test RMSE \\
    \midrule
    무작위 분할 & 0.81 \\
    시간순 분할(뒤 20\% = test) & 44.1 \\
    \bottomrule
  \end{tabular}
  \par\vskip0.5em\normalsize
  \begin{itemize}
    \item \kb{무작위}: test 점들이 시간축 곳곳에 흩어져 각 test 점 근처에
      train 점이 항상 있음(과거든 미래든) → kNN이 잘 맞힘. 그런데
      \textbf{미래 값을 과거 값으로 ``예측''}한 것 — 실전(미래는 아직
      없음)에선 성립하지 않음.
    \item \kb{시간순}: test가 train의 시간 범위(horizon)를 벗어나 kNN이
      밖으로 외삽, 추세를 못 쫓아 RMSE 급등.
  \end{itemize}
  \vskip0.2em
  \begin{block}{교훈}
    이 \textbf{44.1이 실제 일반화 성능에 가까운 값}. 시계열은 무작위가
    아니라 \textbf{시간순으로} 나눠야 하고, 더 근본적으로
    \kq{``실전에서 겹치지 않을 것''을 같은 조각에 두지 말라}.
    {\footnotesize \texttt{train\_test\_split(shuffle=False)} /
    \texttt{TimeSeriesSplit}. 그룹(환자·기기·지점): \texttt{GroupKFold}.}
  \end{block}
\end{frame}
